\section{The Order of Mass: Structure, Retention, Loss, Composition}

The Order of Mass is where the comparison either becomes exact or becomes rhetoric. Both books print it as a continuous rubricated sequence; both can therefore be inventoried unit by unit, and the units can be sorted into four classes---retained with the same words, retained with altered words, dropped, and newly composed---without any interpretive step at all. The sorting is done here first, and interpreted afterwards.

One preliminary matters. The 1962 Order of Mass is printed as a single ceremonial script for a priest at the altar, with variations for solemn Mass supplied in rubrics and in the separate \emph{Ritus servandus}. The postconciliar Order of Mass is printed as one script for Mass with a congregation, followed by a distinct \emph{Ordo Missæ cuius unus tantum minister participat} for Mass with a single server.\endnote{\emph{Missale Romanum}, editio typica tertia, \emph{Ordo Missæ} at printed pp.~503ff.\ (artifact pp.~303ff.), with the shorter order for Mass with one minister at artifact pp.~383ff.; \emph{Missale Romanum} (1962), \emph{Ordo Missæ} at printed p.~216 (artifact p.~297) and \emph{Ritus servandus in celebratione Missæ} at printed pp.~LIV--LXV (artifact pp.~56--67).} The default case has been inverted: what the older book treats as the norm from which solemnity is built up, the newer book treats as the reduction from which the norm is abbreviated. This is a structural fact about the two books, visible from their tables of contents, and it precedes every particular difference.

\subsection{The inventory}

\begin{comparetable}{1962 Order of Mass}{Postconciliar Order of Mass}{Relation}
\emph{In nomine Patris}; antiphon \emph{Introibo ad altare Dei} with Ps 42, said alternately with the ministers, then repeated; \emph{Adiutorium nostrum} & \emph{In nomine Patris}, sung or said with the people, then a greeting (three forms; \emph{Pax vobis} for a bishop) & Psalm and antiphon dropped; the sign of the cross retained and made a common act; a scriptural greeting newly composed from 2 Cor 13:13 and Rom 1:7 alongside \emph{Dominus vobiscum}\\
Priest's \emph{Confiteor}, ministers' \emph{Misereatur}; ministers' \emph{Confiteor}, priest's \emph{Misereatur}; \emph{Indulgentiam}; \emph{Deus tu conversus}, \emph{Ostende nobis}, \emph{Domine exaudi}, \emph{Dominus vobiscum} & Penitential act in one of three forms: \emph{Confiteor} said by all together with \emph{Misereatur nostri}; or the \emph{Miserere nostri} versicles; or the \emph{Qui missus es} invocations with \emph{Kyrie}; sprinkling rite permitted in its place on Sundays & Doubled \emph{Confiteor} collapsed into one said by all; text altered (see below); \emph{Indulgentiam} and the four versicles dropped; two alternative forms and a Sunday substitution newly composed\\
\emph{Aufer a nobis} on ascending; \emph{Oramus te, Domine} at the kiss of the altar & No formula; the rubric prescribes the veneration of the altar by kiss and, as opportune, incensation & Dropped as texts; the gestures retained as rubric\\
Introit antiphon; \emph{Kyrie}; \emph{Gloria} & Entrance chant during the procession; \emph{Kyrie} (unless already used); \emph{Gloria} & Retained; \emph{Gloria} identical in wording; entrance chant relocated to the procession rather than read at the book\\
\emph{Dominus vobiscum}; Collect, one or more as the Office requires & \emph{Oremus}, silence, then one Collect & Retained; the multiplication of collects removed\\
Epistle; Gradual; Alleluia or Tract; Sequence when appointed & First reading; responsorial psalm; second reading on Sundays and solemnities; Alleluia or other chant; Sequence when appointed & Restructured: see section~6\\
\emph{Munda cor meum}; \emph{Iube, Domine, benedicere} (or the deacon's \emph{Iube, domne} with \emph{Dominus sit in corde tuo}) & \emph{Iube, domne, benedicere} and \emph{Dominus sit in corde tuo} when a deacon reads; \emph{Munda cor meum} only when the priest reads without a deacon & Both retained, with the alternatives now assigned by whether a deacon is present rather than by degree of solemnity\\
Gospel; \emph{Per evangelica dicta} said at the kissing of the book & Gospel; \emph{Per evangelica dicta} said quietly at the kissing of the book & Retained verbatim\\
--- & Homily, obligatory on Sundays and holy days of obligation, commended on other days & Newly prescribed within the Order of Mass\\
\emph{Credo} when appointed & \emph{Credo} when appointed; the Apostles' Creed permitted in its place, especially in Lent and Eastertide & Nicene text retained verbatim; a second creed added as an option\\
--- & Universal prayer (prayer of the faithful) & Newly prescribed\\
Offertory antiphon and the offertory prayers & Preparation of the gifts & Treated in section~4\\
Preface; \emph{Sanctus} & Preface; \emph{Sanctus} & Retained; \emph{Sanctus} identical in wording; the repertory of prefaces greatly enlarged\\
Canon Missæ & Prex eucharistica I (Canon romanus) or II, III, IV & Treated in section~5\\
\emph{Pater noster} sung or said by the priest alone, the ministers answering \emph{Sed libera nos a malo} & \emph{Pater noster} said by the priest with the people & Retained; the speaker changed\\
\emph{Libera nos} said silently, with the intercession of Our Lady, Peter, Paul and Andrew; fraction; \emph{Pax Domini}; \emph{Hæc commixtio et consecratio} & \emph{Libera nos} said aloud, without the saints, with \emph{exspectantes beatam spem} added and the people's doxology; \emph{Domine Iesu Christe, qui dixisti} said aloud; \emph{Pax Domini}; optional exchange of peace; fraction; \emph{Hæc commixtio} & Retained with substantial alteration; one of the priest's three private prayers before Communion promoted to an audible prayer for peace; an exchange of peace made an option\\
\emph{Agnus Dei} & \emph{Agnus Dei}, repeatable if the fraction is prolonged & Retained verbatim\\
\emph{Domine Iesu Christe, qui dixisti}; \emph{Domine Iesu Christe, Fili Dei vivi}; \emph{Perceptio Corporis tui} & \emph{Domine Iesu Christe, Fili Dei vivi} \emph{or} \emph{Perceptio Corporis et Sanguinis tui} & All three retained in the book; the first relocated and made audible, the other two reduced to a choice of one\\
\emph{Panem cælestem accipiam}; \emph{Domine, non sum dignus} said three times; \emph{Corpus Domini nostri} & \emph{Ecce Agnus Dei \dots{} Beati qui ad cenam Agni vocati sunt}; \emph{Domine, non sum dignus} said once by priest and people together; \emph{Corpus Christi custodiat me} & \emph{Panem cælestem} dropped; \emph{Ecce Agnus Dei} moved from the Communion of the faithful to this point and expanded with Apoc 19:9; the threefold \emph{Domine, non sum dignus} reduced to one and made common\\
\emph{Quid retribuam}; \emph{Sanguis Domini nostri} & \emph{Sanguis Christi custodiat me} & \emph{Quid retribuam} dropped; the formula shortened\\
Communion of the faithful: \emph{Ecce Agnus Dei}, \emph{Domine non sum dignus} three times, \emph{Corpus Domini nostri Iesu Christi custodiat animam tuam in vitam æternam. Amen} & \emph{Corpus Christi} --- \emph{Amen} & Formula abbreviated and the response transferred to the communicant\\
\emph{Quod ore sumpsimus}; \emph{Corpus tuum, Domine, quod sumpsi} & \emph{Quod ore sumpsimus} & First retained verbatim; second dropped\\
Communion antiphon; Postcommunion(s); \emph{Oratio super populum} on Lenten ferias & Communion chant; silence or a psalm; Prayer after Communion; \emph{Oratio super populum} on Lenten ferias and, as an option, elsewhere & Retained; a period of silence newly prescribed as an option\\
\emph{Ite, missa est} / \emph{Benedicamus Domino} / \emph{Requiescant in pace} & Announcements if needed; greeting; blessing; dismissal & Retained, reordered: in 1962 the dismissal precedes the blessing; in the newer book the blessing precedes the dismissal\\
\emph{Placeat tibi, sancta Trinitas}; blessing; Last Gospel (Jn 1:1--14) & --- & \emph{Placeat} and the Last Gospel dropped\\
\end{comparetable}

\subsection{What the inventory shows}

Four observations follow directly, and none of them requires a judgement.

\textbf{The great fixed texts are untouched.} \emph{Gloria}, \emph{Credo}, \emph{Sanctus}, \emph{Pater noster} and \emph{Agnus Dei} stand in the newer book with the same Latin words they have in the older one. So do \emph{Per evangelica dicta}, \emph{Quod ore sumpsimus}, and---as section~4 will show---\emph{In spiritu humilitatis} and the \emph{Orate, fratres} with its response. Anyone who has been told that the reform rewrote the Ordinary of the Mass is holding a claim that five minutes with the two books disproves.

\textbf{The losses are concentrated at the two ends and in the priest's private apparatus.} What is gone is the psalm and doubled confession at the foot of the altar, the two prayers of the ascent, \emph{Panem cælestem}, \emph{Quid retribuam}, \emph{Corpus tuum Domine}, \emph{Placeat tibi}, and the Last Gospel; together with the offertory prayers treated separately below. Every one of these except the Last Gospel is a text the older book directs the priest to say \emph{secreto} or in a low voice, and most are in the first person singular. The reform did not remove the congregation's texts; it removed a large part of the celebrant's private ones.

\textbf{Where a text was retained but altered, the alteration is usually an addition or a reassignment, not a subtraction.} The \emph{Confiteor} gains \emph{et omissione} to the list of ways one has sinned---in thought, word, deed \emph{and omission}---and loses the enumeration of Michael, John the Baptist, Peter and Paul in favour of ``all the Angels and Saints''; it is said once, by everyone, rather than twice, by two parties in turn.\endnote{1962: ``\dots{} quia peccavi nimis cogitatione, verbo et opere \dots{} Ideo precor beatam Mariam semper Virginem, beatum Michaelem Archangelum, beatum Ioannem Baptistam, sanctos Apostolos Petrum et Paulum, omnes Sanctos, et vos, fratres \dots{}'' (\emph{Missale Romanum} 1962, \emph{Ordo Missæ} n.~1017, artifact p.~297, collated visually 25 July 2026). 2002: ``\dots{} quia peccavi nimis cogitatione, verbo, opere et omissione \dots{} Ideo precor beatam Mariam semper Virginem, omnes Angelos et Sanctos, et vos, fratres \dots{}'' (\emph{Ordo Missæ} n.~4, artifact p.~303).} The \emph{Libera nos} loses ``past, present and to come'' and the four named intercessors and gains \emph{exspectantes beatam spem et adventum Salvatoris nostri Iesu Christi} from Titus 2:13; it also moves from silence to audible prayer and acquires a people's doxology.\endnote{1962: \emph{Ordo Missæ} n.~1117, artifact p.~396: ``Libera nos, quæsumus Domine, ab omnibus malis, præteritis, præsentibus et futuris: et intercedente beata et gloriosa semper Virgine Dei Genetrice Maria, cum beatis Apostolis tuis Petro et Paulo, atque Andrea, et omnibus Sanctis, \dots{} da propitius pacem in diebus nostris \dots{}'' said with the paten, silently. 2002: \emph{Ordo Missæ} n.~125, artifact p.~365, said \emph{manibus extensis} and concluded by the people's \emph{Quia tuum est regnum}.} A reader who wants to argue that these changes impoverish the rite has a real argument available; a reader who wants to argue that they suppress something is misdescribing the text.

\textbf{The newly composed material is small in bulk and large in position.} What is genuinely new in the Order of Mass---as distinct from relocated or shortened---amounts to the alternative greetings, two of the three penitential forms, the sprinkling option, the memorial acclamation, the people's doxology after \emph{Libera nos}, the \emph{Beati qui ad cenam Agni vocati sunt}, and the two \emph{Benedictus es} prayers at the preparation of the gifts. That is a short list. But almost every item on it is placed at a seam where the people speak, and their cumulative effect is to convert a rite in which the congregation had four or five spoken moments into one in which it has a dozen. Whether that is the recovery of something ancient or the addition of something modern is precisely the argument; that it is the reform's structural centre of gravity is not in dispute, and the 1969 constitution says as much when it explains that the Order was revised so that ``the parts proper to each, whether minister or people, might be more clearly distinguished.''

\subsection{Options: the difference the inventory cannot show}

An inventory of units understates the structural difference between the two books, because it counts what is printed and not how many of the printed things are alternatives.

The 1962 Order of Mass is, with a few defined exceptions, an invariable script. The exceptions are real but narrow: the psalm \emph{Iudica me} and the \emph{Gloria Patri} are omitted in Masses for the dead and from Passion Sunday to Maundy Thursday; the Gloria and Creed are said when the Office requires; the number of collects and postcommunions follows the calendar; the incensations belong to solemn Mass; the dismissal takes one of three forms; the Last Gospel is proper on Palm Sunday and omitted in defined cases. Beyond these, the celebrant chooses nothing. He does not choose a greeting, a penitential formula, a creed, a Eucharistic Prayer, a preface where one is appointed, an acclamation, or a Communion prayer. The book tells him what to say.

The postconciliar Order of Mass is built on choice. At the opening there are three greetings. At the penitential act there are three forms and, on Sundays, a fourth possibility in the sprinkling rite. There are two creeds. There are four Eucharistic Prayers in the Order itself, and further ones in an appendix---the third typical edition's promulgating decree names ``the Eucharistic Prayer for Reconciliation'' and ``a particular Eucharistic Prayer which may be used for various needs''---and, for most days, a repertory of prefaces from which the priest selects.\endnote{Congregation for Divine Worship and the Discipline of the Sacraments, decree \emph{De editione typica tertia}, Prot.\ N.~143/00/L (20 April 2000), artifact p.~3: ``In appendice ad Ordinem Missæ inveniuntur etiam Prex Eucharistica pro Reconciliatione, necnon Prex Eucharistica peculiaris, quæ pro variis necessitatibus adhiberi potest.''} After the consecration there are three memorial acclamations. Before Communion there are two prayers of which the priest says one. The sign of peace is offered ``pro opportunitate.'' A period of silence after Communion is one of three permitted uses of that moment. And article~31 of the \emph{Institutio generalis} allows the priest, where the rubrics provide for it, to adapt the prescribed introductions ``somewhat, so that they respond to the understanding of those participating''---while charging him to keep the sense of what the Missal proposes and to express it in few words.\endnote{\emph{Institutio Generalis Missalis Romani} (2002), n.~31, artifact p.~25: ``Ubi a rubricis statuitur, celebranti licet eas aliquatenus aptare ut participantium captui respondeant; curet tamen sacerdos ut sensum monitionis quæ in Missali proponitur ipse semper servet eamque paucis verbis exprimat.'' The article also reserves to him the introductions before the Mass of the day, the liturgy of the word, and the Eucharistic Prayer---the last ``numquam vero intra Precem ipsam,'' never within the Prayer itself.}

This is the difference that no side-by-side table registers, and it is arguably the most consequential of all. A book of alternatives produces variation between celebrations of the same rite on the same day. That variation may be experienced as pastoral responsiveness or as instability, and the reform's defenders and critics have both understood its significance: \emph{Institutio generalis} n.~31 hedges the adaptation it permits with a threefold restraint, and the concern the hedges answer is precisely the concern the critics raise. What can be said flatly is that the older book locates the rite's fixity in its text, and the newer book locates it in a framework within which texts are selected. That is a change of kind, not of degree, and it happened at the level of structure rather than of wording.

\subsection{The claim that does not survive: ``the 1962 book is the unreformed rite''}

The comparison above is conventionally staged as ``the traditional Mass'' against ``the reformed Mass.'' The 1962 book will not support the first half of that description, and it says so itself in several places.

Its own opening decree, quoted in section~1, presents it as an edition prepared to conform to a code of rubrics approved in July 1960. That code, printed in the same volume as the \emph{Rubricæ generales Missalis romani}, had already done substantial work. It reduced the octaves of the entire Roman calendar to three---``celebrantur tantum octavæ Nativitatis Domini, Paschatis et Pentecostes, exclusis omnibus aliis''---suppressing every other octave in the universal and every particular calendar.\endnote{\emph{Missale Romanum} (1962), \emph{Rubricæ generales Missalis romani} n.~64, printed p.~XV (artifact p.~13).} It replaced the old ranks of doubles and semidoubles with four classes of liturgical day. It ruled that in solemn Mass the celebrant ``omits those things which are proclaimed by the sacred ministers or by a lector''---abolishing the duplication by which the priest had read privately what the subdeacon and deacon were singing.\endnote{\emph{Rubricæ generales} n.~513(f), printed p.~XXXV (artifact p.~33), collated visually 25 July 2026: ``In Missa solemni, celebrans: \dots{} \emph{f}) omittit ea quæ a ministris sacris vel a lectore proferuntur.'' The same number's preceding paragraphs assign him what he does sing and say; n.~514 extends the rule, with adjustments, to sung Mass without sacred ministers. The corresponding provision in \emph{Ritus servandus} VI, 4 has the subdeacon sing the Epistle ``quam celebrans sedens auscultat''---which the celebrant listens to, seated.} It permitted the celebrant at solemn and sung Mass to sit during \emph{Kyrie}, \emph{Gloria}, sequence and \emph{Credo}.\endnote{\emph{Rubricæ generales} n.~523, printed p.~XXXVI (artifact p.~34).} And it suppressed the \emph{Confiteor} and absolution before the Communion of the faithful within Mass, retaining only \emph{Ecce Agnus Dei} and the threefold \emph{Domine, non sum dignus}.\endnote{\emph{Rubricæ generales} n.~503, printed p.~XXXIV (artifact p.~32): ``Quoties sancta Communio infra Missam distribuitur, celebrans, sumpto sacratissimo Sanguine, omissis confessione et absolutione, dictis tamen \emph{Ecce Agnus Dei} et ter \emph{Domine, non sum dignus}, immediate ad distributionem sanctæ Eucharistiæ procedit.''}

These are not marginal adjustments. Three of them---the end of celebrant duplication, the celebrant's seat, the suppressed second \emph{Confiteor}---are moves in exactly the direction the postconciliar reform would continue, made by the very book that is now offered as the alternative to it. The honest description of the comparison is therefore not ``unreformed against reformed'' but ``the reform of 1960 against the reform of 1969.'' That description costs the critical case nothing that matters; it only removes an argument the critical case does not need and cannot sustain.
